\documentclass[a4paper,10pt]{report}\input{../0.Préambule/Préambule_cours_prof}\begin{document}

\pagestyle{empty}
\newpage
\section*{Produits de nombres relatifs \hspace{3cm} Sujet A}
\addcontentsline{toc}{section}{Produits de nombres relatifs}

\noindent \textit{L'usage de la calculatrice est \textbf{interdite}.} \hfill \textit{Durée : 30 minutes}

\paragraph{Exercice 1}: Question de cours\hfill \textbf{2 points}

Donner votre définition de l'opposé d'un nombre, ainsi qu'un exemple.

\paragraph{Exercice 2}: \hfill \textbf{6 points}

Effectue les calculs suivants (on détaillera toutes les étapes du calcul):
\begin{enumerate}
\begin{tabularx}{\linewidth}{*{2}{>{\arraybackslash}X}}
\item $S=(-3,5)\times (-10)\times (+20)\times (-0,3)$&
\item $A=(-75)\times (-0,25)\times (+4)\times (+2)$\\
\item $C=(-1,5)\times (+6)\times (-1)\times (+0,7)\times (-10)$&\\
\end{tabularx}
\end{enumerate}

\paragraph{Exercice 3}: \hfill \textbf{6 points}

Effectue les calculs suivants (on détaillera toutes les étapes du calcul):
\begin{enumerate}
\begin{tabularx}{\linewidth}{*{3}{>{\arraybackslash}X}}
\item $L=-3+(1-5)\times (-6)$&
\item $E=5-3\times 3+2\times (-5)$&
\item $O=1+(-2)\times (-2)-(-3)\times (-3)$\\
\end{tabularx}
\end{enumerate}

\paragraph{Exercice 4}: \hspace{6.5cm} \textbf{6 points}

On considère le programme de calcul suivant :

\vspace{-10mm}
\begin{flushright}
\arrayrulecolor{black}
\begin{tabular}{|p{6cm}|}
\hline
Choisir un nombre\\
Le multiplier par $(-3)$\\
Ajouter le double du nombre de départ\\\hline
\end{tabular}
\end{flushright}

\vspace{-10mm}
\begin{enumerate}
\item Quel résultat obtient-on lorsque le nombre de départ est :
\begin{enumerate}
\begin{tabularx}{\linewidth}{*{3}{>{\arraybackslash}X}}
\item $7$ ?&
\item $-5$ ?&
\item $3,6$ ? \\
\end{tabularx}
\end{enumerate}
\item Que remarque-t-on ?
\end{enumerate}

\vspace{3cm}
\section*{Produits de nombres relatifs \hspace{3cm} Sujet B}
\addcontentsline{toc}{section}{Produits de nombres relatifs}

\noindent \textit{L'usage de la calculatrice est \textbf{interdite}.} \hfill \textit{Durée : 30 minutes}

\paragraph{Exercice 1}: Question de cours\hfill \textbf{2 points}

Donner votre définition de l'opposé d'un nombre, ainsi qu'un exemple.

\paragraph{Exercice 2}: \hfill \textbf{6 points}

Effectue les calculs suivants (on détaillera toutes les étapes du calcul):
\begin{enumerate}
\begin{tabularx}{\linewidth}{*{2}{>{\arraybackslash}X}}
\item $S=(-3,2)\times (-10)\times (+2)\times (-0,5)$&
\item $A=(-100)\times (-0,25)\times (+3)\times (+2)$\\
\item $C=(-1,5)\times (+4)\times (-1)\times (+0,8)\times (-3)$&\\
\end{tabularx}
\end{enumerate}

\paragraph{Exercice 3}: \hfill \textbf{6 points}

Effectue les calculs suivants (on détaillera toutes les étapes du calcul):
\begin{enumerate}
\begin{tabularx}{\linewidth}{*{3}{>{\arraybackslash}X}}
\item $L=-2+(1-4)\times (-7)$&
\item $E=1-2\times 3+4\times (-5)$&
\item $O=2+(-2)\times (-3)-(-4)\times (-3)$\\
\end{tabularx}
\end{enumerate}

\paragraph{Exercice 4}: \hspace{6.5cm} \textbf{6 points}

On considère le programme de calcul suivant :

\vspace{-10mm}
\begin{flushright}
\arrayrulecolor{black}
\begin{tabular}{|p{6cm}|}
\hline
Choisir un nombre\\
Le multiplier par $(-3)$\\
Ajouter le double du nombre de départ\\\hline
\end{tabular}
\end{flushright}

\vspace{-10mm}
\begin{enumerate}
\item Quel résultat obtient-on lorsque le nombre de départ est :
\begin{enumerate}
\begin{tabularx}{\linewidth}{*{3}{>{\arraybackslash}X}}
\item $6$ ?&
\item $-8$ ?&
\item $2,4$ ? \\
\end{tabularx}
\end{enumerate}
\item Que remarque-t-on ?
\end{enumerate}
\end{document}
